\section{Problem Formulation}\label{sec:method:formulation}

\subsection*{Observed input}

The data consist of $S$ \emph{recordings} indexed by $s \in \{1,\ldots,S\}$, each an intercepted stream in which pulses from several emitters are interleaved in time.
Recording $s$ is observed as a finite, time-ordered sequence of \glspl{pdw}
\[
  \mat{X}^{(s)} = (\vect{x}^{(s)}_1,\ldots,\vect{x}^{(s)}_{N_s}),
  \qquad
  \vect{x}^{(s)}_n \in \R^{d},
\]
where each $\vect{x}^{(s)}_n$ contains the measured parameters of one pulse (carrier frequency, amplitude, pulse width, time of arrival, and angles of arrival), and the length $N_s$ varies between recordings.
The encoder operates on fixed-length windows extracted from these recordings as described in \cref{sec:method:data:window}.
When considering one recording or window, the superscript is omitted.

\subsection*{Labels and target partitions}

Each pulse $n \in \{1,\ldots,N_s\}$ carries an emitter identity $e^{(s)}_n \in \{1,\ldots,K^{(s)}\}$ and an operating mode $m^{(s)}_n \in \{1,\ldots,M\}$.
The emitter count $K^{(s)}$ may vary between recordings, and each emitter identifier is meaningful only within its recording.
Mode labels, by contrast, are drawn from a finite catalogue of size $M$ shared across all recordings.
A recording generally contains only $M^{(s)} \leq M$ of these modes.
The two labels describe different relations: one emitter may switch between several modes, and the same mode may occur on several emitters.
The two labels of a pulse are collected into the ground-truth label pair
\begin{equation}\label{eq:method:label_pair}
  \vect{y}^{(s)}_n = \bigl(e^{(s)}_n,\, m^{(s)}_n\bigr)
  \in \{1,\ldots,K^{(s)}\} \times \{1,\ldots,M\},
\end{equation}
which constitutes the complete per-pulse annotation of recording $s$.

For recording $s$, the coordinates of $\vect{y}^{(s)}_n$ induce emitter and mode cells on the pulse-index set,
\begin{equation}\label{eq:method:emitter_partition}
  \mathcal{E}^{(s)}_k = \{\, n : e^{(s)}_n = k \,\},
  \qquad
  k \in \{1,\ldots,K^{(s)}\},
\end{equation}
\begin{equation}\label{eq:method:mode_partition}
  \mathcal{M}^{(s)}_\ell = \{\, n : m^{(s)}_n = \ell \,\},
  \qquad
  \ell \in \{1,\ldots,M\}.
\end{equation}
The emitter cells are non-empty and partition $\{1,\ldots,N_s\}$.
The mode cells are also disjoint, but $\mathcal{M}^{(s)}_\ell$ is empty when mode $\ell$ is absent from the recording; the non-empty mode cells form the corresponding partition.
Recovering the emitter partition is the deinterleaving task, whereas recovering the mode partition groups pulses by the active mode for each PDW.

\begin{figure}[tbp]
  \centering
  % =====================================================================
% Schematic of the two label-induced partitions of a single pulse train
% used in \cref{sec:method:formulation}.  Each \gls{pdw} carries two
% labels: an emitter identity (encoded by fill colour) and an operating
% mode (encoded by border style: solid, dashed, dotted).  These two
% labellings induce two distinct partitions of the same index set:
% the emitter partition $\{\mathcal{E}_k\}$ (within-train) and the mode
% partition $\{\mathcal{M}_\ell\}$ (defined relative to a global mode
% catalogue).  Loaded from
% sections/04_method/01_problem_formulation.tex via \input{...}.
% =====================================================================
\begin{tikzpicture}[
    font=\small,
    >=Latex,
    pulse/.style={
      circle, draw,
      minimum size=8mm,
      inner sep=0pt,
      line width=0.55pt,
    },
    em1/.style={fill=blue!22},
    em2/.style={fill=red!22},
    em3/.style={fill=green!32},
    modeA/.style={},
    modeB/.style={densely dashed, line width=0.85pt},
    modeC/.style={densely dotted, line width=1.1pt},
    brace/.style={inner xsep=1pt, inner ysep=0pt},
    rowlabel/.style={font=\footnotesize, anchor=east, align=right},
  ]

  \def\dx{0.95}

  % ------------------------------------------------------------------
  % Top row: input pulse train.
  % Per-pulse labels used throughout the figure:
  %   x1: em1,A   x2: em2,A   x3: em1,B   x4: em3,A   x5: em2,B
  %   x6: em3,C   x7: em1,C   x8: em1,A   x9: em2,C
  % ------------------------------------------------------------------
  \node[pulse, em1, modeA] (t1) at (0*\dx, 0) {$\vect{x}_1$};
  \node[pulse, em2, modeA] (t2) at (1*\dx, 0) {$\vect{x}_2$};
  \node[pulse, em1, modeB] (t3) at (2*\dx, 0) {$\vect{x}_3$};
  \node[pulse, em3, modeA] (t4) at (3*\dx, 0) {$\vect{x}_4$};
  \node[pulse, em2, modeB] (t5) at (4*\dx, 0) {$\vect{x}_5$};
  \node[pulse, em3, modeC] (t6) at (5*\dx, 0) {$\vect{x}_6$};
  \node[pulse, em1, modeC] (t7) at (6*\dx, 0) {$\vect{x}_7$};
  \node[pulse, em1, modeA] (t8) at (7*\dx, 0) {$\vect{x}_8$};
  \node[pulse, em2, modeC] (t9) at (8*\dx, 0) {$\vect{x}_9$};

  \node[brace, anchor=east] (top_lb) at ([xshift=-1mm]t1.west) {$\Bigg\{$};
  \node[brace, anchor=west] (top_rb) at ([xshift=1mm]t9.east)  {$\Bigg\}$};

  % ------------------------------------------------------------------
  % Forking arrows from the input to each partition.
  % ------------------------------------------------------------------
  \coordinate (fork)   at (4*\dx, -0.7);
  \coordinate (em_in)  at (4*\dx, -1.7);
  \coordinate (md_in)  at (4*\dx, -3.7);

  \draw[thick] (4*\dx, -0.55) -- (fork);
  \draw[thick, ->] (fork) -- (em_in)
        node[midway, right, font=\footnotesize]
            {by emitter $e_n$};
  \draw[thick, ->] (fork) -- (md_in)
        node[midway, right, font=\footnotesize]
            {by mode $m_n$};

  % ------------------------------------------------------------------
  % Emitter partition (within-train).
  %   E1 (blue):  x1, x3, x7, x8
  %   E2 (red):   x2, x5, x9
  %   E3 (green): x4, x6
  % ------------------------------------------------------------------
  \def\yem{-2.5}
  \node[pulse, em1, modeA] (e1) at (0.0*\dx, \yem) {$\vect{x}_1$};
  \node[pulse, em1, modeB] (e3) at (1.0*\dx, \yem) {$\vect{x}_3$};
  \node[pulse, em1, modeC] (e7) at (2.0*\dx, \yem) {$\vect{x}_7$};
  \node[pulse, em1, modeA] (e8) at (3.0*\dx, \yem) {$\vect{x}_8$};

  \node[pulse, em2, modeA] (e2) at (4.3*\dx, \yem) {$\vect{x}_2$};
  \node[pulse, em2, modeB] (e5) at (5.3*\dx, \yem) {$\vect{x}_5$};
  \node[pulse, em2, modeC] (e9) at (6.3*\dx, \yem) {$\vect{x}_9$};

  \node[pulse, em3, modeA] (e4) at (7.6*\dx, \yem) {$\vect{x}_4$};
  \node[pulse, em3, modeC] (e6) at (8.6*\dx, \yem) {$\vect{x}_6$};

  \node[brace, anchor=east] (em1lb) at ([xshift=-0.5mm]e1.west) {$\Big\{$};
  \node[brace, anchor=west] (em1rb) at ([xshift=0.5mm]e8.east)  {$\Big\},$};
  \node[brace, anchor=east] (em2lb) at ([xshift=-0.5mm]e2.west) {$\Big\{$};
  \node[brace, anchor=west] (em2rb) at ([xshift=0.5mm]e9.east)  {$\Big\},$};
  \node[brace, anchor=east] (em3lb) at ([xshift=-0.5mm]e4.west) {$\Big\{$};
  \node[brace, anchor=west] (em3rb) at ([xshift=0.5mm]e6.east)  {$\Big\}$};

  \node[brace, anchor=east] (em_lb) at ([xshift=-1.5mm]em1lb.west) {$\Bigg\{$};
  \node[brace, anchor=west] (em_rb) at ([xshift=1.5mm]em3rb.east)  {$\Bigg\}$};

  \node[rowlabel] at ([xshift=-3mm]em_lb.west)
        {emitter partition\\$\{\mathcal{E}_k\}_{k=1}^{K_{\mathrm{em}}}$};

  % ------------------------------------------------------------------
  % Mode partition (global catalogue).
  %   A (solid):  x1, x2, x4, x8
  %   B (dashed): x3, x5
  %   C (dotted): x6, x7, x9
  % ------------------------------------------------------------------
  \def\ymd{-4.5}
  \node[pulse, em1, modeA] (m1) at (0.0*\dx, \ymd) {$\vect{x}_1$};
  \node[pulse, em2, modeA] (m2) at (1.0*\dx, \ymd) {$\vect{x}_2$};
  \node[pulse, em3, modeA] (m4) at (2.0*\dx, \ymd) {$\vect{x}_4$};
  \node[pulse, em1, modeA] (m8) at (3.0*\dx, \ymd) {$\vect{x}_8$};

  \node[pulse, em1, modeB] (m3) at (4.3*\dx, \ymd) {$\vect{x}_3$};
  \node[pulse, em2, modeB] (m5) at (5.3*\dx, \ymd) {$\vect{x}_5$};

  \node[pulse, em3, modeC] (m6) at (6.6*\dx, \ymd) {$\vect{x}_6$};
  \node[pulse, em1, modeC] (m7) at (7.6*\dx, \ymd) {$\vect{x}_7$};
  \node[pulse, em2, modeC] (m9) at (8.6*\dx, \ymd) {$\vect{x}_9$};

  \node[brace, anchor=east] (md1lb) at ([xshift=-0.5mm]m1.west) {$\Big\{$};
  \node[brace, anchor=west] (md1rb) at ([xshift=0.5mm]m8.east)  {$\Big\},$};
  \node[brace, anchor=east] (md2lb) at ([xshift=-0.5mm]m3.west) {$\Big\{$};
  \node[brace, anchor=west] (md2rb) at ([xshift=0.5mm]m5.east)  {$\Big\},$};
  \node[brace, anchor=east] (md3lb) at ([xshift=-0.5mm]m6.west) {$\Big\{$};
  \node[brace, anchor=west] (md3rb) at ([xshift=0.5mm]m9.east)  {$\Big\}$};

  \node[brace, anchor=east] (md_lb) at ([xshift=-1.5mm]md1lb.west) {$\Bigg\{$};
  \node[brace, anchor=west] (md_rb) at ([xshift=1.5mm]md3rb.east)  {$\Bigg\}$};

  \node[rowlabel] at ([xshift=-3mm]md_lb.west)
        {mode partition\\$\{\mathcal{M}_\ell\}_{\ell=1}^{M}$};

\end{tikzpicture}

  \caption{Each pulse $\vect{x}_n$ is assigned both an emitter label $e_n$ (color) and a mode label $m_n$ (border style).}
  \label{fig:method:partitions}
\end{figure}

\subsection*{Estimation objective}

At inference, only the window $\mat{X}$ is observed; the label pairs in \cref{eq:method:label_pair} and the numbers of groups present are unknown.
The task is to estimate for every pulse a label pair
\begin{equation}\label{eq:method:estimated_label_pair}
  \hat{\vect{y}}_n = \bigl(\hat{e}_n,\, \hat{m}_n\bigr)
  \in \{1,\ldots,\hat{K}\} \times \{1,\ldots,\hat{M}\},
  \qquad n \in \{1,\ldots,L\},
\end{equation}
where the cardinalities $\hat{K}$ and $\hat{M}$ must themselves be estimated from the window.
The coordinates of these estimates induce the estimated partitions
\begin{equation}\label{eq:method:estimated_partitions}
  \widehat{\mathcal{E}}_k
  =
  \{\,n : \hat{e}_n = k\,\},
  \quad k \in \{1,\ldots,\hat{K}\},
  \qquad
  \widehat{\mathcal{M}}_\ell
  =
  \{\,n : \hat{m}_n = \ell\,\},
  \quad \ell \in \{1,\ldots,\hat{M}\}.
\end{equation}
The estimate $\hat{\vect{y}}_n$ approximates $\vect{y}_n$ not elementwise but through these partitions: the objective is to recover, within each window, the restrictions of both ground-truth partitions in \cref{eq:method:emitter_partition,eq:method:mode_partition} to that window, up to independent permutations of the cluster indices of each coordinate.
The estimated symbols themselves carry no meaning; agreement depends only on which pulses are grouped together, and is quantified by the permutation-invariant partition-similarity metrics defined in \cref{sec:method:metrics}.

Both coordinates of $\hat{\vect{y}}_n$ are produced per window, so their symbols are local to that window.
Emitter identities are additionally meaningful only within the recording from which the window was drawn.
Neither $\hat{K}$ nor $\hat{M}$ is assumed known, and $\hat{M}$ is not constrained by the training catalogue size $M$; this permits clusters corresponding to modulation types not observed during training.

The proposed method realizes this estimation with a shared encoder and two embedding branches, one per coordinate of $\hat{\vect{y}}_n$: recording-local emitter labels supervise one branch, while globally shared mode labels make the other branch comparable across recordings.
The supervised contrastive losses are specified in \cref{sec:method:loss}, and the architecture, together with the post-hoc step that turns embeddings into the discrete estimates, is described in \cref{sec:method:arch}.
